\documentclass[14pt,t,dvipsnames]{beamer}
% \setbeamerfont{alerted text}{series=\bfseries\large}
\setbeamertemplate{navigation symbols}{}
\usefonttheme[onlymath]{serif}

\setlength{\parskip}{1em}


\usepackage[utf8]{inputenc}
\usepackage[T1]{fontenc}
\usepackage{lmodern}
\usepackage{makecell} %...used in tabular for \thead
 
%...For highlighting
\usepackage{xcolor, soul} 
\makeatletter
\let\HL\hl
\renewcommand\hl{%
  \let\set@color\beamerorig@set@color
  \let\reset@color\beamerorig@reset@color
  \HL}
\makeatother

\hypersetup{
    colorlinks=true,
    linkcolor=blue,
    filecolor=magenta,      
    urlcolor=cyan,
    pdftitle={Overleaf Example},
    pdfpagemode=FullScreen,
}

\usepackage[acronym]{glossaries}
\glsdisablehyper %...disables hyperlinks for glossary terms in frames


\usepackage{amsmath}  %...for the equation* environment

%...Tikz and PGF packages
\usepackage{tikz}
\usetikzlibrary{calc,math}
\usetikzlibrary{patterns}
\usepackage{tikz-cd}
\usepackage{pgfplots}
\pgfplotsset{compat=1.18}
\usetikzlibrary{backgrounds} %...For shading portions of graphs
% \usetikzlibrary{math}
\tikzset{>=latex} %...For pretty arrows
\tikzset{
    declare function = {trajectoryequationTEN(\x,\vi,\thetai)= tan(\thetai)*\x - 10*\x^2/(2*(\vi*cos(\thetai))^2);},
    declare function = {trajectoryequation(\x,\vi,\thetai)= tan(\thetai)*\x - 9.8*\x^2/(2*(\vi*cos(\thetai))^2);},
    declare function = {patheq(\x,\yi,\vi,\thetai)= \yi + tan(\thetai)*\x - 9.8*\x^2/(2*(\vi*cos(\thetai))^2);}
} %...Plots the trajectory equation for projectile motion:


\usetikzlibrary{positioning}
\tikzset{onslide/.code args={<#1>#2}{%
  \only<#1>{\pgfkeysalso{#2}} % \pgfkeysalso doesn't change the path
}}

%...SI Units
\usepackage{siunitx}
%\sisetup{per-mode=fraction}
\usepackage{cancel}
\DeclareSIUnit{\mile}{mi}


%...Symbols packages
\usepackage{fontawesome} % for \faHome
\usepackage{tikzsymbols} % for tree and stick figure
\usepackage{twemojis} %...for twitter emojis
\usepackage{utfsym}  %...for a basic star symbol

%...Breaks long URLs
\usepackage{xurl}

%...Multiple columns for table of contents
\usepackage{multicol}

%...beamer specific
\AtBeginSection[]
{
  \begin{frame}
    \frametitle{Table of Contents}
    \begin{multicols}{2}
        \tableofcontents[currentsection]
    \end{multicols}
  \end{frame}
}

\setbeamerfont{frametitle}{size=\normalsize}

%...Vector commands
\newcommand{\bvec}[1]{\vec{\mathbf{#1}}}
\newcommand{\bhat}[1]{\,\hat{\mathbf{#1}}}


%%%%%%% DST, XVT, and VAT Triangles
\newcommand{\DSTtriangle}[1][]{%
    % Key-value setup for the shade argument
    \pgfkeys{/DST/.cd, shade/.store in=\DSTshade, shade=none, #1}%
    
    \begin{tikzpicture}[x=1.5cm,y=1.5cm]
        \coordinate (D) at (0,1);
        \coordinate (S) at ({cos(210)},{sin(210)});
        \coordinate (T) at ({cos(330)},{sin(330)});
        
        % Midpoints
        \coordinate (M1) at ($(D)!0.5!(S)$);
        \coordinate (M2) at ($(D)!0.5!(T)$);
        \coordinate (M3) at ($(M1)!0.5!(M2)$);
        \coordinate (M4) at ($(S)!0.5!(T)$);
        
        % Variables for comparison
        \def\shadeD{D}
        \def\shadeS{S}
        \def\shadeT{T}
        
        % Shading logic (drawn first so lines appear on top)
        \ifx\DSTshade\shadeD
            \fill[gray!30] (D) -- (M1) -- (M2) -- cycle;
        \fi
        \ifx\DSTshade\shadeS
            \fill[gray!30] (S) -- (M4) -- (M3) -- (M1) -- cycle;
        \fi
        \ifx\DSTshade\shadeT
            \fill[gray!30] (T) -- (M4) -- (M3) -- (M2) -- cycle;
        \fi
        
        % Triangle outer border
        \draw (D) -- (S) -- (T) -- cycle;
        
        % Lines between midpoints
        \draw (M1) -- (M2) node[pos=0.5,above=0.03] {D};
        \draw (M3) -- (M4) node[pos=0.5,left=0.1] {S} node[pos=0.5,right=0.06] {T};
    \end{tikzpicture}
}

\newcommand{\XVTtriangle}[1][]{%
    % Key-value setup for the shade argument
    \pgfkeys{/DST/.cd, shade/.store in=\DSTshade, shade=none, #1}%
    
    \begin{tikzpicture}[x=1.5cm,y=1.5cm]
        \coordinate (D) at (0,1);
        \coordinate (S) at ({cos(210)},{sin(210)});
        \coordinate (T) at ({cos(330)},{sin(330)});
        
        % Midpoints
        \coordinate (M1) at ($(D)!0.5!(S)$);
        \coordinate (M2) at ($(D)!0.5!(T)$);
        \coordinate (M3) at ($(M1)!0.5!(M2)$);
        \coordinate (M4) at ($(S)!0.5!(T)$);
        
        % Variables for comparison
        \def\shadeD{X}
        \def\shadeS{V}
        \def\shadeT{T}
        
        % Shading logic (drawn first so lines appear on top)
        \ifx\DSTshade\shadeD
            \fill[gray!30] (D) -- (M1) -- (M2) -- cycle;
        \fi
        \ifx\DSTshade\shadeS
            \fill[gray!30] (S) -- (M4) -- (M3) -- (M1) -- cycle;
        \fi
        \ifx\DSTshade\shadeT
            \fill[gray!30] (T) -- (M4) -- (M3) -- (M2) -- cycle;
        \fi
        
        % Triangle outer border
        \draw (D) -- (S) -- (T) -- cycle;
        
        % Lines between midpoints
        \draw (M1) -- (M2) node[pos=0.5,above=0.03] {$\Delta x$};
        \draw (M3) -- (M4) node[pos=0.5,shift={(-0.3,-0.01)}] {$\bar{v}$} node[pos=0.5,shift={(0.3,-0.01)}] {$\Delta t$};
    \end{tikzpicture}
}

\newcommand{\VATtriangle}[1][]{%
    % Key-value setup for the shade argument
    \pgfkeys{/DST/.cd, shade/.store in=\DSTshade, shade=none, #1}%
    
    \begin{tikzpicture}[x=1.5cm,y=1.5cm]
        \coordinate (D) at (0,1);
        \coordinate (S) at ({cos(210)},{sin(210)});
        \coordinate (T) at ({cos(330)},{sin(330)});
        
        % Midpoints
        \coordinate (M1) at ($(D)!0.5!(S)$);
        \coordinate (M2) at ($(D)!0.5!(T)$);
        \coordinate (M3) at ($(M1)!0.5!(M2)$);
        \coordinate (M4) at ($(S)!0.5!(T)$);
        
        % Variables for comparison
        \def\shadeD{V}
        \def\shadeS{A}
        \def\shadeT{T}
        
        % Shading logic (drawn first so lines appear on top)
        \ifx\DSTshade\shadeD
            \fill[gray!30] (D) -- (M1) -- (M2) -- cycle;
        \fi
        \ifx\DSTshade\shadeS
            \fill[gray!30] (S) -- (M4) -- (M3) -- (M1) -- cycle;
        \fi
        \ifx\DSTshade\shadeT
            \fill[gray!30] (T) -- (M4) -- (M3) -- (M2) -- cycle;
        \fi
        
        % Triangle outer border
        \draw (D) -- (S) -- (T) -- cycle;
        
        % Lines between midpoints
        \draw (M1) -- (M2) node[pos=0.5,above=0.03] {$\Delta v$};
        \draw (M3) -- (M4) node[pos=0.5,shift={(-0.3,-0.01)}] {$\bar{a}$} node[pos=0.5,shift={(0.3,-0.01)}] {$\Delta t$};
    \end{tikzpicture}
}

%...Speedometer
\newcommand{\speedometer}[3]{%
    % #1: X-shift (y is fixed at 2) (e.g., 1)
    % #2: Needle rotation angle (e.g., 20)
    % #3: Top label (e.g., \texttt{Speed})
    % #4: Bottom readout (e.g., \textsf{32.0\,m/s})
    \begin{scope}[shift={#1},x=1.3cm,y=1.3cm]
        \draw (0,0) circle (1);
        \foreach \x in {340,360,...,560} {
            \coordinate (P) at ({cos(\x)},{sin(\x)});
            \draw[ultra thick,gray] (P) -- ($(P)!0.16!(0,0)$);
        }
        \foreach \x in {350,370,...,550} {
            \coordinate (P) at ({cos(\x)},{sin(\x)});
            \draw[thick,gray] (P) -- ($(P)!0.08!(0,0)$);
        }
        \fill (0,0) circle (0.02);
        \draw[fill=lightgray,rotate={#2},red] (0,0.02) rectangle (0.975,-0.02);
        \node at (0,0.4) {\texttt{Speed}};
        \draw[rounded corners] (-0.65,-0.75) rectangle ++(1.3,0.5) node[pos=0.5] {\textsf{#3\,m/s}};
    \end{scope}
}


\usepackage[acronym]{glossaries}
\usepackage{makecell}

\setbeamerfont{alerted text}{series=\bfseries}

\glsdisablehyper

\makeglossaries
\input{2027-physics/glossary}
%...UNIT 1: CONSTANT VELOCITY

\newcommand{\xygraphI}[3]{
\begin{tikzpicture}
    \begin{axis}[height=7.7cm,width=8.5cm,
        axis lines=center,
        ylabel={$y$},
        xlabel={$x$},
        ymin=-10,ymax=10,
        xmin=-10,xmax=10,
        ytick={-10,-8,...,10},
        xtick={-10,-8,...,10},
        grid=both, 
        minor y tick num=1,
        minor x tick num=1,
    ]
        \addplot[thick,domain=-10:10] {#1*x + #2};
        \ifprintanswers
            \addplot[red,only marks,mark=*] coordinates{#3};
        \fi
    \end{axis}
\end{tikzpicture}
}


\newcommand{\distancegraphI}[1]{
    \begin{tikzpicture}
        \begin{axis}[height=6.5cm,width=6.8cm,
            axis lines=left,
            xlabel = {Time (hours)},
            ylabel = {Distance (miles)},
            ymin=0, ymax=100,
            xmin=0, xmax=5,
            ytick={0,10,...,100},
            xtick={0,1,...,5},
            minor x tick num=1,
            grid=both,
            ]
            \addplot[very thick,domain=0:5] {#1*x};
        \end{axis}
    \end{tikzpicture}
}

\newcommand{\positiongraphI}[3]{
    \begin{tikzpicture}
        \begin{axis}[height=5.6cm,width=6.5cm,
            axis lines=left,
            xlabel = {Time (s)},
            ylabel = {Position (m)},
            ymin=0, ymax=20,
            xmin=0, xmax=14,
            ytick={0,4,...,20},
            xtick={0,2,...,14},
            minor y tick num=3,
            minor x tick num=1,
            grid=both,
            ]
            \addplot[very thick,domain=0:14] {#1*x + #2};
            \ifprintanswers
                \addplot[red,only marks,mark=*] coordinates{#3};
            \fi
        \end{axis}
    \end{tikzpicture}

%...AI-CHATBOT PROMPT:
% In an xy plane, the y limits are 0 to 20, and the x limits are 0 to 14. Give me 5 excellent linear equations (y = mx + b) for students to practice slope calculations. Output the results in the form \positiongraphI{m}{b}{coordinates} where coordinates should be in the form (x1,y1)(x2,y2), representing 2 points on the line for calculating slope. 
}

\newcommand{\positionfigureI}[1]{
\begin{center}
    \begin{tikzpicture}[y=1cm,x=1cm]
        \foreach \x in {-5,-4,...,5} {
            \draw[ultra thin] (\x,-1mm) -- (\x,1mm) node[below=2mm] {$\x$};
        }
        \draw[->] (-5,0) -- (5,0) node[below=6mm,pos=0.5] {Position (m)};
        \node[above] at (#1,0) {\Strichmaxerl[2]};
        \ifprintanswers
            \draw[very thick,->,red] (0,0.7) -- ++(#1,0) node[above,pos=0.5] {$\vec{x}$};
        \fi
    \end{tikzpicture}  
\end{center}
}

\newcommand{\positionfigureII}[2]{
\begin{center}
\begin{tikzpicture}[y=1cm,x=1cm]
    \foreach \x in {-5,-4,...,5} {
        \draw[ultra thin] (\x,-1mm) -- (\x,1mm) node[below=3mm] {$\x$};
    }
    \draw[->] (-5,0) -- (5,0) node[below=7mm,pos=0.5] {Position (m)};
    \fill (#1,0) circle (1mm) node[above=1mm] {A};
    \fill (#2,0) circle (1mm) node[above=1mm] {B};
\end{tikzpicture}
\end{center}
}


\newcommand{\positionfigureIII}[4]{
\begin{center}
\begin{tikzpicture}[y=3mm,x=8mm]
    \draw[black!10,xstep=1,ystep=2] (0,0) grid (12,4);
    \foreach \x in {0,1,...,12} {
        \draw[ultra thin] (\x,-1mm) -- (\x,1mm) node[below=3mm] {$\x$};
    }
    \draw[->] (0,0) -- (12,0) node[right=1mm] {$x$ (m)};
    \draw[ultra thick,->,rounded corners=0.7mm] (#1,1)
        -- (#2,1) 
        -- (#2,2) 
        -- (#3,2) 
        -- (#3,3) 
        -- (#4,3);
    \node[fill=black!10] at (#1,-3) {$x_i$};
    \node[above=1mm,fill=black!10] at (#4,4) {$x_f$};
\end{tikzpicture}
\end{center}
}

\newcommand{\positionfigureIV}[2]{
    \begin{center}
    \begin{tikzpicture}[y=1.4mm,x=1.4mm]
        \foreach \x in {-50,-40,...,50} {
            \draw[ultra thin] (\x,-1mm) -- (\x,1mm) node[below=3mm] {$\x$};
        }
        \foreach \x in {-45,-35,...,45} {
            \draw[ultra thin] (\x,-0.8mm) -- (\x,0.8mm);
        }
        \draw[->] (-50,0) -- (50,0) node[below=7mm,pos=0.5] {Position (m)};
        \fill (#1,0) circle (1mm) node[above=1mm] {A};
        \fill (#2,0) circle (1mm) node[above=1mm] {B};
    \end{tikzpicture}
    \end{center}
}

\newcommand{\positionfigureV}[6]{
    %...PRO TIP: Type the following into ChatGPT:
    % Generate 5 integers between -10 to 10 (inclusive) such that the absolute value of the sum of the 5 integers is less than or equal to 10. Print each of the integers in each of the first 5 curly brackets below. Calculate the maximum displacement from equilibrium and, if it's greater than 10, enter it in the 6th curly bracket; otherwise enter 10 for the 6th curly bracket. Do not output in markdown; I just want the verbatim text: \positionfigureV{}{}{}{}{}{}
    
    \begin{tikzpicture}[x=4mm,y=4mm]
        \draw[step=1,lightgray] (-10,0) grid (10,4);
        \draw[->,line width=0.6mm,-{Triangle}] (0,4) -- ++(#5,0);
        \draw[->,line width=0.6mm,-{Triangle}] (0,3) -- ++(#4,0);
        \draw[->,line width=0.6mm,-{Triangle}] (0,2) -- ++(#3,0);
        \draw[->,line width=0.6mm,-{Triangle}] (0,1) -- ++(#2,0);
        \draw[->,line width=0.6mm,-{Triangle}] (0,0) -- ++(#1,0);
        \node[below] at (current bounding box.south) {Figure 1: Displacement vectors.};
    \end{tikzpicture}
    
    \smallskip
    
    \begin{tikzpicture}[x=4mm,y=4mm]
        \draw[step=1,lightgray] (-#6,0) grid (#6,4);
        \draw (0,0) node {\scriptsize $|$} node[below=2pt] {$0$};
            \begin{scope}[red,opacity={\ifprintanswers 1\else 0\fi}]
                \draw[->,line width=0.3mm,-{Triangle}] ({#1+#2+#3+#4},4) -- ++(#5,0);
                \draw[->,line width=0.3mm,-{Triangle}] ({#1+#2+#3},3) -- ++(#4,0);
                \draw[->,line width=0.3mm,-{Triangle}] ({#1+#2},2) -- ++(#3,0);
                \draw[->,line width=0.3mm,-{Triangle}] (#1,1) -- ++(#2,0);
                \draw[->,line width=0.3mm,-{Triangle}] (0,0) -- ++(#1,0);
                \pgfmathsetmacro{\vecsum}{#1+#2+#3+#4+#5}
        
                \pgfmathtruncatemacro{\iszero}{(\vecsum==0)}
                
                \ifnum\iszero=1
                    \node[above] at (0,4.6) {vector sum is zero};
                \else
                    \draw[->,line width=0.6mm,-{Triangle}] (0,4.8) -- ++(\vecsum,0)
                        node[above=1mm,pos=0.5] {vector sum};
                \fi
        \end{scope}
        \node[below] at (current bounding box.south) {Figure 2: Head-to-tail addition and vector sum.};
    \end{tikzpicture}
}

\newcommand{\positionfigureVI}[1]{
\begin{tikzpicture}[y=6mm,x=6mm]
    \draw[black!20,step=1] (0,0) grid (24,2);  
    \draw[|-|] (0,-0.3) -- ++(1,0) node[below] {\small 1 meter};
    \ifprintanswers
    \foreach \x in {0,#1,...,24}{
        \fill[red] (\x,1) circle (0.9mm);
    }
    \fi
\end{tikzpicture}
}

\newcommand{\positiongraphIII}[6]{
    %...Generates a quadrant I position vs. time with linear function.
    % #1 = velocity (slope)
    % #2 = initial position (y-intercept)
\begin{tikzpicture}
\begin{axis}[height=5.9cm,width=6.4cm,
    axis lines=left,
    xlabel = {Time (s)},
    ylabel = {Position (m/s)},
    ymin=0, ymax=10,
    xmin=0, xmax=10,
    ytick={0,1,...,10},
    xtick={0,1,...,10},
    grid = both,
    clip=true,
    legend style={at={(1.1,0.85)}}
    % set layers
    ]
    \addplot[very thick,black,domain=0:10] {#1*\x + #2};
    \addlegendentry{A} % <-- Legend for the first plot
    
    \addplot[densely dashed,very thick,black,domain=0:10] {#3*\x + #4};
    \addlegendentry{B} 

    \addplot[only marks, mark=*,mark size=1.5pt,black,domain=0:10] {#5*\x + #6};
    \addlegendentry{C} 
    \end{axis}
\end{tikzpicture}
}

\newcommand{\velocitygraph}[2]{
    \begin{tikzpicture}
        \begin{axis}[height=6cm,width=7cm,
            axis y line=left,
            axis x line=center,
            xlabel = {Time (s)},
            ylabel = {Velocity (m/s)},
            ymin=-5, ymax=5,
            xmin=0, xmax=10,
            ytick={-5,-4,...,5},
            xtick={0,2,...,10},
            grid=both,
            minor x tick num = 1,
            x label style={at={(axis description cs:1,0.5)},right=1mm},
            clip=false
        ]
            \addplot[very thick,black,domain=0:#2] {#1};
            \ifprintanswers
                \begin{pgfonlayer}{background}
                    \fill[red!12] (0,0) rectangle (#2,#1);
                \end{pgfonlayer}
            \fi
            \end{axis}
    \end{tikzpicture}
}

%...AI-CHATBOT PROMPT:
%Generate a random integer between -5 to 5 (inclusive), and another one between 5 and 10 (inclusive). Output the results in the form \velocitygraph{a}{b} where a and b are the generated numbers.




%..UNIT 2: ACCELERATION

\newcommand{\velocitytableI}[4]{%
    % #1 = initial velocity
    % #2 = acceleration
    % #3 = initial time
    % #4 = change in time
\begin{tabular}{|c|c|}
    \hline
    \thead{\bfseries Velocity (m/s)} & \thead{\bfseries Time (s)} \\ \hline
    $\fpeval{#1 + #2*(#3 + 0*#4)}$ & \fpeval{#3 + 0*#4} \\ \hline
    $\fpeval{#1 + #2*(#3 + 1*#4)}$ & \fpeval{#3 + 1*#4} \\ \hline
    $\fpeval{#1 + #2*(#3 + 2*#4)}$ & \fpeval{#3 + 2*#4} \\ \hline
    $\fpeval{#1 + #2*(#3 + 3*#4)}$ & \fpeval{#3 + 3*#4} \\ \hline
    $\fpeval{#1 + #2*(#3 + 4*#4)}$ & \fpeval{#3 + 4*#4} \\ \hline
\end{tabular}
}



\newcommand{\positiontableI}[5]{%
    % #1 = initial position
    % #2 = initial velocity
    % #3 = initial time
    % #4 = change in time
    % #5 = acceleration
\begin{tabular}{|c|c|}
    \hline
    \thead{\bfseries Position (m)} & \thead{\bfseries Time (s)} \\ \hline
    $\fpeval{#1 + #2*(#3 + 0*#4) + 0.5*#5*(#3 + 0*#4)^2}$ & \fpeval{#3 + 0*#4} \\ \hline
    $\fpeval{#1 + #2*(#3 + 1*#4) + 0.5*#5*(#3 + 1*#4)^2}$ & \fpeval{#3 + 1*#4} \\ \hline
    $\fpeval{#1 + #2*(#3 + 2*#4) + 0.5*#5*(#3 + 2*#4)^2}$ & \fpeval{#3 + 2*#4} \\ \hline
    $\fpeval{#1 + #2*(#3 + 3*#4) + 0.5*#5*(#3 + 3*#4)^2}$ & \fpeval{#3 + 3*#4} \\ \hline
    $\fpeval{#1 + #2*(#3 + 4*#4) + 0.5*#5*(#3 + 4*#4)^2}$ & \fpeval{#3 + 4*#4} \\ \hline
\end{tabular}
}

\NewDocumentCommand{\velocitygraphI}{m m o}{
    % Produces a line on a velocity vs time graph.
    % #1 = slope (acceleration) 
    % #2 = y-intercept (initial velocity)
    % #3 [optional] = sample coordinates (x1,y1)(x2,y2) for finding slope
\begin{tikzpicture}
\begin{axis}[height=6cm,width=6.4cm,
    axis y line=left,
    axis x line=center,
    ylabel = {Velocity (m/s)},
    xlabel = {Time (s)},
    ymin=-8, ymax=8,
    xmin=0, xmax=10,
    ytick={-8,-6,...,8},
    xtick={0,1,...,10},
    grid = both,
    minor y tick num=1,
    ]
    \addplot[very thick,black,domain=0:10] {#1*\x + #2};
    \ifprintanswers
            % \IfValueT checks if the 3rd argument was actually provided
            \IfValueT{#3}{
                \addplot[red,only marks,mark=*] coordinates{#3};
            }
        \fi
    \end{axis}
\end{tikzpicture}

%...AI CHATBOT PROMPT:
% A velocity vs. time graph has y limits -8 to 8 and x limits 0 to 10. Generate a linear function for students to practice calculating acceleration. Output the result in the form \velocitygraphI{m}{b}{(x1,y1)(x2,y2)} where m is the slope (acceleration), b is the y-intercept (initial velocity), and (x1,y1)(x2,y2) are two coordinates on the line used to calculate slope. Negative accelerations are acceptable.
}

\newcommand{\positiongraphII}[3]{
    % Generates a position vs. time graph for uniform acceleration motion.
    %...Inspired by https://ophysics.com/k4.html
    % #1 = initial position
    % #2 = initial velocity
    % #3 = acceleration
\begin{tikzpicture}
    \begin{axis}[height=6cm,width=5cm,
        axis y line=left,
        axis x line=center,
        ylabel={Position (m)},
        xlabel={Time (s)},
        ymin=-4,ymax=4,
        xmin=0,xmax=4,
        ytick={-4,-3,...,4},
        xtick={0,1,...,4},
        grid=both,
        x label style={at={(axis description cs: 1,0.5)},right=1mm},
    ]
        \addplot[very thick,domain=0:4] {#1 + #2*x + 0.5*#3*x^2};        
    \end{axis}
\end{tikzpicture}    
}

\newcommand{\velocitygraphII}[2]{
    %...Generates a minimalistic v vs. t graph.
    % #1 = initial velocity (y-intercept)
    % #2 = acceleration (slope)
    %... See https://ophysics.com/k4.html
\begin{tikzpicture}[y=5mm,x=7mm]
    \draw[->] (0,-4.4) -- (0,4.4) node[left,pos=0.9] {$v$} node[pos=0.5,left] {\footnotesize $0$};
    \draw[->] (0,0) -- (4.4,0) node[right] {$t$};
    \draw[very thick,domain=0:4] plot(\x,{#1 + #2*\x});
\end{tikzpicture}
}


\newcommand{\velocitygraphIII}[3]{
    %...Generates a quadrant I velocity vs. time with linear function.
    % #1 = acceleration (slope)
    % #2 = initial velocity (y-intercept)
    % #3 = final time (where line stops)
\begin{tikzpicture}
\begin{axis}[height=6cm,width=6.4cm,
    axis lines=left,
    xlabel = {Time (s)},
    ylabel = {Velocity (m/s)},
    ymin=0, ymax=10,
    xmin=0, xmax=10,
    ytick={0,1,...,10},
    xtick={0,1,...,10},
    grid = both,
    clip=true,
    % set layers
    ]
    \addplot[very thick,black,domain=0:#3] {#1*\x + #2};
    \ifprintanswers
    \begin{pgfonlayer}{background}
        \fill[red!10] (0,#2) -- (#3,{#1*#3+#2}) -- (#3,0) -- (0,0) -- cycle;
    \end{pgfonlayer}
    \fi
    \end{axis}
\end{tikzpicture}

%...AI CHATBOT PROMPT:
% A velocity vs. time graph has y limits 0 to 10 and x limits 0 to 10. Generate a linear function for students to practice calculating displacement (as area under the curve). Output the result in the form \velocitygraphIII{m}{b} where m is the slope (acceleration) and b is the y-intercept (initial velocity). Negative accelerations are acceptable.
}

\newcommand{\triangleI}[3]{%
    % #1 = base
    % #2 = height
    % #3 = a parameter to control figure size
    \begin{tikzpicture}[y=1mm*#3,x=1mm*#3]
        \draw[thick,fill=black!10] (-0.5*#1,0) -- (0.5*#1,0) -- (0,#2) -- cycle;
        \draw[densely dashed] (0,0) -- (0,#2) node[pos=0.4,right] {#2};
        \draw[|-|] (-0.5*#1,-2mm) -- (0.5*#1,-2mm) node[below,pos=0.5] {#1};
    \end{tikzpicture}
}

\newcommand{\triangleII}[3]{%
    % #1 = base
    % #2 = height
    % #3 = a parameter to control figure size
    \begin{tikzpicture}[y=1mm*#3,x=1mm*#3]
        \draw[thick,fill=black!10] (0,0) -- (#1,0) node[below,pos=0.5] {#1} -- (#1,#2) node[right,pos=0.5] {#2} -- cycle;
    \end{tikzpicture}
}


% \renewcommand*{\glstextformat}[1]{\textcolor{black}{#1}} %...without this, glossary terms are colored blue in Beamer, BUT not needed if using \glsdisablehyper. Consider deleting this line.

% \newcommand{\DSTtriangle}{
% \begin{tikzpicture}[x=1.5cm,y=1.5cm]
%     \coordinate (D) at (0,1);
%     \coordinate (S) at ({cos(210)},{sin(210)});
%     \coordinate (T) at ({cos(330)},{sin(330)});
%     \draw (D) -- (S) -- (T) -- (D);
    
%     % Midpoints
%     \coordinate (M1) at ($(D)!0.5!(S)$);
%     \coordinate (M2) at ($(D)!0.5!(T)$);
%     \coordinate (M3) at ($(M1)!0.5!(M2)$);
%     \coordinate (M4) at ($(S)!0.5!(T)$);
    
%     % Lines between midpoints
%     \draw (M1) -- (M2) node[pos=0.5,above=0.03] {D};
%     \draw (M3) -- (M4) node[pos=0.5,left=0.1] {S} node[pos=0.5,right=0.06] {T};
% \end{tikzpicture}
% }

\newcommand{\DSTtriangle}[1][]{%
    % Key-value setup for the shade argument
    \pgfkeys{/DST/.cd, shade/.store in=\DSTshade, shade=none, #1}%
    
    \begin{tikzpicture}[x=1.5cm,y=1.5cm]
        \coordinate (D) at (0,1);
        \coordinate (S) at ({cos(210)},{sin(210)});
        \coordinate (T) at ({cos(330)},{sin(330)});
        
        % Midpoints
        \coordinate (M1) at ($(D)!0.5!(S)$);
        \coordinate (M2) at ($(D)!0.5!(T)$);
        \coordinate (M3) at ($(M1)!0.5!(M2)$);
        \coordinate (M4) at ($(S)!0.5!(T)$);
        
        % Variables for comparison
        \def\shadeD{D}
        \def\shadeS{S}
        \def\shadeT{T}
        
        % Shading logic (drawn first so lines appear on top)
        \ifx\DSTshade\shadeD
            \fill[gray!30] (D) -- (M1) -- (M2) -- cycle;
        \fi
        \ifx\DSTshade\shadeS
            \fill[gray!30] (S) -- (M4) -- (M3) -- (M1) -- cycle;
        \fi
        \ifx\DSTshade\shadeT
            \fill[gray!30] (T) -- (M4) -- (M3) -- (M2) -- cycle;
        \fi
        
        % Triangle outer border
        \draw (D) -- (S) -- (T) -- cycle;
        
        % Lines between midpoints
        \draw (M1) -- (M2) node[pos=0.5,above=0.03] {D};
        \draw (M3) -- (M4) node[pos=0.5,left=0.1] {S} node[pos=0.5,right=0.06] {T};
    \end{tikzpicture}
}

\title{Unit 1: Constant Velocity}
\subtitle{Physics}
\author{}
\date{Updated on \today}

\AtBeginSection[]
{
  \begin{frame}{Table of Contents}
        \tableofcontents[currentsection]
  \end{frame}
}


\begin{document}

\begin{frame}
\titlepage
\end{frame}

\section{Mon Aug 17 - position, distance, displacement}

\begin{frame}[label=5i5c0]
\textbf{\gls{number line}}: \glsdesc{number line} 

\begin{center}
\begin{tikzpicture}[y=1cm,x=1cm]
    \foreach \x in {-5,-4,...,5} {
        \draw[ultra thin] (\x,-1mm) -- (\x,1mm) node[below=2mm] {$\x$};
    }
    \draw[->] (-5,0) -- (5,0);
\end{tikzpicture} 
\end{center}

\pause


\textbf{\gls{position}}: \glsdesc{position}
\vspace{-1ex}

\begin{center}
\begin{tikzpicture}[y=1cm,x=1cm]
    \foreach \x in {0,1,...,10} {
        \draw[ultra thin] (\x,-1mm) -- (\x,1mm) node[below=2mm] {$\x$};
    }
    \draw[->] (0,0) -- (10,0) node[below=6mm,pos=0.5] {Position (m)};
    \node[above] at (3,0) {\Strichmaxerl[2]};
\end{tikzpicture} 
\end{center}

\end{frame}

\begin{frame}

\textbf{\gls{motion}}: \glsdesc{motion} 

\medskip

\textbf{\gls{one-dimensional motion}}: \glsdesc{one-dimensional motion} 

\begin{center}
\begin{tikzpicture}[y=1cm,x=1cm]
    \foreach \x in {0,1,...,10} {
        \draw[ultra thin] (\x,-1mm) -- (\x,1mm) node[below=2mm] {$\x$};
    }
    \draw[->] (0,0) -- (10,0) node[below=6mm,pos=0.5] {Position (m)};
    \foreach \x in {3, 5, 7, 9}{
        \node[above] at (\x,0) {\Strichmaxerl[2]};
    }
\end{tikzpicture} 
\end{center}


\end{frame}

\begin{frame}
\textbf{\gls{distance}} ($d$): \glsdesc{distance}    

\begin{center}
\begin{tikzpicture}[y=3mm,x=8.5mm]
    \draw[black!10,xstep=1,ystep=2] (0,0) grid (10,4);
    \foreach \x in {0,1,...,10} {
        \draw[ultra thin] (\x,-1mm) -- (\x,1mm) node[below=3mm] {$\x$};
    }
    \draw[->] (0,0) -- (10,0) node[right=1mm] {$x$ (m)};
    \draw[ultra thick,->,rounded corners=0.7mm] (0,1)
        -- (5,1) 
        -- (5,2) 
        -- (2,2) 
        -- (2,3) 
        -- (8,3);
    \node[fill=black!10] at (0,-4) {$x_i$};
    \node[above=1mm,fill=black!10] at (8,4) {$x_f$};
\end{tikzpicture}
\end{center}
\vspace{-1em}

\pause

\begin{flalign*}
    x_i &= \text{initial position (starting point)} && \\[1ex]
    x_f &= \text{final position (ending point)} && \\[1ex]
\end{flalign*}
\end{frame}

\begin{frame}
\textit{Example.} What is the distance traveled?

\begin{center}
\begin{tikzpicture}[y=3mm,x=8.5mm]
    \draw[black!10,xstep=1,ystep=2] (0,0) grid (10,4);
    \foreach \x in {0,1,...,10} {
        \draw[ultra thin] (\x,-1mm) -- (\x,1mm) node[below=3mm] {$\x$};
    }
    \draw[->] (0,0) -- (10,0) node[right=1mm] {$x$ (m)};
    \draw[ultra thick,->,rounded corners=0.7mm] (0,1)
        -- (5,1) 
        -- (5,2) 
        -- (2,2) 
        -- (2,3) 
        -- (8,3);
    \node[fill=black!10] at (0,-4) {$x_i$};
    \node[above=1mm,fill=black!10] at (8,4) {$x_f$};
\end{tikzpicture}
\end{center}

\pause

\begin{equation*}
    d = \SI{5}{m} + \SI{3}{m} + \SI{6}{m} = \boxed{\SI{14}{m}}
\end{equation*}
\end{frame}

\begin{frame}
\textbf{\gls{displacement}} ($\Delta x$): \glsdesc{displacement}    

\begin{center}
\begin{tikzpicture}[y=3mm,x=8.5mm]
    \draw[black!10,xstep=1,ystep=2] (0,0) grid (10,4);
    \foreach \x in {0,1,...,10} {
        \draw[ultra thin] (\x,-1mm) -- (\x,1mm) node[below=3mm] {$\x$};
    }
    \draw[->] (0,0) -- (10,0) node[right=1mm] {$x$ (m)};
    \draw[ultra thick,->,rounded corners=0.7mm] (0,1)
        -- (5,1) 
        -- (5,2) 
        -- (2,2) 
        -- (2,3) 
        -- (8,3);
    \node[fill=black!10] at (0,-4) {$x_i$};
    \node[above=1mm,fill=black!10] at (8,4) {$x_f$};
\end{tikzpicture}
\end{center}

\vspace{-1cm}

\begin{flalign*}
    x_i &= \text{initial position (starting point)} && \\[1ex]
    x_f &= \text{final position (ending point)} && \\[1ex]
    \Delta x &= x_f - x_i = \text{final pos.} - \text{initial pos.}
\end{flalign*}
\end{frame}

\begin{frame}
\textit{Example.} What is the object's displacement?   

\begin{center}
\begin{tikzpicture}[y=3mm,x=8.5mm]
    \draw[black!10,xstep=1,ystep=2] (0,0) grid (10,4);
    \foreach \x in {0,1,...,10} {
        \draw[ultra thin] (\x,-1mm) -- (\x,1mm) node[below=3mm] {$\x$};
    }
    \draw[->] (0,0) -- (10,0) node[right=1mm] {$x$ (m)};
    \draw[ultra thick,->,rounded corners=0.7mm] (0,1)
        -- (5,1) 
        -- (5,2) 
        -- (2,2) 
        -- (2,3) 
        -- (8,3);
    \node[fill=black!10] at (0,-4) {$x_i$};
    \node[above=1mm,fill=black!10] at (8,4) {$x_f$};
\end{tikzpicture}
\end{center}

\vspace{-1cm}

\pause

\begin{align*}
    \Delta x &= x_f - x_i \\[1ex]
    &= \SI{8}{m} - \SI{0}{m} \\[1ex]
    &= \boxed{+\SI{8}{m}}
\end{align*}
\end{frame}

\begin{frame}
An object's position in one-dimensional motion is represented by\ldots

\begin{enumerate}[A]
    \item \alert<2->{a number line}
    \item a curvy path
    \item an equation
    \item a unit circle
\end{enumerate}

\uncover<2->{
\begin{tikzpicture}[y=1cm,x=1cm]
    \foreach \x in {0,1,...,10} {
        \draw[ultra thin] (\x,-1mm) -- (\x,1mm) node[below=2mm] {$\x$};
    }
    \draw[->] (0,0) -- (10,0) node[below=6mm,pos=0.5] {Position (m)};
    \node[above] at (3,0) {\Strichmaxerl[2]};
\end{tikzpicture} 
}
\end{frame}

\begin{frame}
Distance means\ldots 

\begin{enumerate}[A]
    \item the net change in an object's position.
    \item \alert<2->{the length of a path traveled.}
    \item how fast an object is moving.
    \item the direction in which an object moves.
\end{enumerate}
\end{frame}

\begin{frame}
True or false? Distance and displacement mean the same thing.

\begin{enumerate}[A]
    \item True
    \item \alert<2->{False}
\end{enumerate}

\vspace{-1.5cm}

\uncover<2->{
\begin{center}
\huge
\begin{equation*}
    \mathrm{distance \neq displacement}
\end{equation*}
\end{center}
}
\end{frame}

\begin{frame}
Find the object's displacement.

\begin{center}
\begin{tikzpicture}[y=3mm,x=8.5mm]
    \draw[black!10,xstep=1,ystep=2] (0,0) grid (10,4);
    \foreach \x in {0,1,...,10} {
        \draw[ultra thin] (\x,-1mm) -- (\x,1mm) node[below=3mm] {$\x$};
    }
    \draw[->] (0,0) -- (10,0) node[right=1mm] {$x$ (m)};
    \draw[ultra thick,->,rounded corners=0.7mm] (7,1) node[left] {\small $x_i$}
        -- (9,1) 
        -- (9,2) 
        -- (1,2) 
        -- (1,3) 
        -- (3,3) node[right] {\small $x_f$};
\end{tikzpicture}
\end{center}

\begin{enumerate}[A]
    \item $+\SI{4}{m}$
    \item \alert<2->{$-\SI{4}{m}$}
    \item $+\SI{12}{m}$
    \item $-\SI{4}{m}$
\end{enumerate}
\end{frame}

\begin{frame}
What distance did the object travel?

\begin{center}
\begin{tikzpicture}[y=3mm,x=8.5mm]
    \draw[black!10,xstep=1,ystep=2] (0,0) grid (10,4);
    \foreach \x in {0,1,...,10} {
        \draw[ultra thin] (\x,-1mm) -- (\x,1mm) node[below=3mm] {$\x$};
    }
    \draw[->] (0,0) -- (10,0) node[right=1mm] {$x$ (m)};
    \draw[ultra thick,->,rounded corners=0.7mm] (7,1) node[left] {\small $x_i$}
        -- (9,1) 
        -- (9,2) 
        -- (1,2) 
        -- (1,3) 
        -- (3,3) node[right] {\small $x_f$};
\end{tikzpicture}
\end{center}

\begin{enumerate}[A]
    \item $+\SI{4}{m}$
    \item $-\SI{4}{m}$
    \item \alert<2->{$+\SI{12}{m}$}
    \item $-\SI{4}{m}$
\end{enumerate}
\end{frame}

\section{Tue Aug 18 - speed, average speed}

\begin{frame}{Warm-Up}
    \begin{enumerate}
        \item In your own words, write the definition of \textbf{position}. 
        \item Draw a visual representation of an object's position.
    \end{enumerate}
\end{frame}

\begin{frame}
\textbf{\gls{speed}}: \uncover<2->{\glsdesc{speed}}

\medskip

\textbf{\gls{average speed}}: \uncover<3->{\glsdesc{average speed}}

\medskip


\uncover<4->{\Large
\begin{equation*}
    \mathrm{speed} = \mathrm{\frac{distance}{time}}
\end{equation*}
}
\end{frame}

\begin{frame}
\textit{Example.} An object travels 8 meters in 4 seconds. 

\begin{tikzpicture}[y=1cm,x=1cm]
    \foreach \x in {0,1,...,10} {
        \draw[ultra thin] (\x,-1mm) -- (\x,1mm) node[below=2mm] {$\x$};
    }
    \draw[->] (0,0) -- (10,0) node[below=6mm,pos=0.5] {Position (m)};
    \node[above] at (8,0) {\Strichmaxerl[2]};
\end{tikzpicture} 

What is its average speed?

\pause

\begin{equation*}
    \mathrm{speed} = \mathrm{\frac{distance}{time}} = \frac{\SI{8}{m}}{\SI{4}{s}} = \boxed{\SI{2}{m/s}}
\end{equation*}
\vspace{1ex}

\pause

The average speed is 2 meters per second.
\end{frame}

\begin{frame}
\textit{Example.} An object travels 10 meters in 2 seconds. 

\begin{tikzpicture}[y=1cm,x=1cm]
    \foreach \x in {0,1,...,10} {
        \draw[ultra thin] (\x,-1mm) -- (\x,1mm) node[below=2mm] {$\x$};
    }
    \draw[->] (0,0) -- (10,0) node[below=6mm,pos=0.5] {Position (m)};
    \node[above] at (0,0) {\Strichmaxerl[2]};
\end{tikzpicture} 

What is its average speed?

\pause

\begin{equation*}
    \mathrm{speed} = \mathrm{\frac{distance}{time}} = \frac{\SI{10}{m}}{\SI{2}{s}} = \boxed{\SI{5}{m/s}}
\end{equation*}
\vspace{1ex}

\pause

The average speed is 2 meters per second.
\end{frame}

\begin{frame}
Use the \textbf{DST triangle} for word problems involving distance, speed, and time.

\begin{center}
\DSTtriangle % No shading
\end{center}

\begin{center}
    D = distance \qquad 
    S = speed \qquad
    T = time
\end{center}
\end{frame}

\begin{frame}
\centering

\fbox{
\begin{minipage}{4.5cm}
    \centering
    \DSTtriangle[shade=S]  

    \vspace{-8mm}

    \begin{equation*}
        \mathrm{speed} = \mathrm{\frac{distance}{time}}
    \end{equation*}
\end{minipage}}%
\hspace{5mm}
\uncover<2->{\fbox{
\begin{minipage}{4.5cm}
    \centering
    \DSTtriangle[shade=T]  

    \vspace{-8mm}

    \begin{equation*}
        \mathrm{time} = \mathrm{\frac{distance}{speed}}
    \end{equation*}
\end{minipage}
}}

\vspace{1em}

\uncover<3->{\fbox{
\begin{minipage}{5.4cm}
    \centering
    \DSTtriangle[shade=D]  

    \vspace{-8mm}

    \begin{equation*}
        \mathrm{distance} = \mathrm{speed \times time}
    \end{equation*}
\end{minipage}
}}
\end{frame}



\begin{frame}
\textit{Example.} An object travels at a speed of 7 meters per second. How how far does it travel after 3 seconds?
\vspace{-5mm}

\pause

\begin{center}
\DSTtriangle[shade=D]  
\end{center}

\vspace{-1cm}

\begin{align*}
    \uncover<3->{\mathrm{distance} &= \mathrm{speed} \times \mathrm{time} \\[1ex]}
    \uncover<4->{&= \SI{7}{m/s} \times \SI{3}{s} \\[1ex]}
    \uncover<5->{&= \boxed{\SI{21}{m}}}
\end{align*}

\uncover<6->{
The object travels a distance of 21 meters.
}
\end{frame}

\begin{frame}
\textit{CFU.} An object travels at a speed of 8 meters per second. How how far does it travel after 4 seconds?

\pause

\begin{enumerate}[A]
    \item \SI{12}{m}
    \item \SI{4}{m}
    \item \alert<2->{\SI{32}{m}}
    \item \SI{24}{m}
\end{enumerate}

\pause

% \begin{center}
% \DSTtriangle[shade=D]  
% \end{center}

\vspace{-1cm}

\begin{align*}
    \mathrm{distance} &= \mathrm{speed} \times \mathrm{time} \\[1ex]
    &= \SI{8}{m/s} \times \SI{4}{s} \\[1ex]
    &= \boxed{\SI{32}{m}}
\end{align*}

% \uncover<6->{
% The object travels a distance of 21 meters.
% }
\end{frame}

\section{Wed Aug 19 - slope, distance vs time graph}

\begin{frame}
\vspace{-5mm}

\begin{minipage}[t]{3.9cm}
    \vspace{0pt}
    \textit{Example.}\\Find the slope of the line.

    \vspace{-5mm}
    \uncover<2->{
    \begin{align*}
        m &= \mathrm{\frac{rise}{run}} \\[1ex]
        &= \frac{y_2 - y_1}{x_2 - x_1} \\[1ex]
        &= \frac{2-0}{1-0}\\[1ex]
        &= 2
    \end{align*}}
\end{minipage}%
\begin{minipage}[t]{6.1cm}
\vspace{0pt}
\centering
\small
\begin{tikzpicture}
    \begin{axis}[height=6.5cm,width=7cm,
        axis lines=center,
        ylabel={$y$},
        xlabel={$x$},
        ymin=-4,ymax=4,
        xmin=-4,xmax=4,
        ytick={-4,-3,...,4},
        xtick={-4,-3,...,4},
        grid=both, 
    ]
        \addplot[very thick,domain=-4:4] {2*x};
        \addplot[red,only marks,mark=*] coordinates{(0,0)(1,2)};
        \draw[red,->] (0,0) --  (0.8,-1) node[below] {$(0,0)$};
        \node[red,right=1mm] at (1,2) {$(1,2)$};
    \end{axis}
\end{tikzpicture}
\end{minipage}

\begin{center}
\uncover<3->{The slope of the line is 2.}
\end{center}
\end{frame}

\begin{frame}
\vspace{-5mm}

\begin{minipage}[t]{3.9cm}
    \vspace{0pt}
    \textit{Example.}\\Find the slope of the line.

    \vspace{-5mm}

    \begin{align*}
        m &= \mathrm{\frac{rise}{run}} \\[1ex]
        &= \frac{y_2 - y_1}{x_2 - x_1} \\[1ex]
        &= \frac{1-3}{3-0}\\[1ex]
        &= -\frac{2}{3}
    \end{align*}
\end{minipage}%
\begin{minipage}[t]{6.1cm}
\vspace{0pt}
\centering
\small
\begin{tikzpicture}
    \begin{axis}[height=6.5cm,width=7cm,
        axis lines=center,
        ylabel={$y$},
        xlabel={$x$},
        ymin=-4,ymax=4,
        xmin=-4,xmax=4,
        ytick={-4,-3,...,4},
        xtick={-4,-3,...,4},
        grid=both, 
    ]
        \addplot[very thick,domain=-4:4] {-2/3*x + 3};
        \addplot[red,only marks,mark=*] coordinates{(0,3)(3,1)};
        % \draw[red,->] (0,0) --  (0.8,-1) node[below] {$(0,0)$};
        \node[red,right=2mm] at (0,3) {$(0,3)$};
        \node[red,above=1mm] at (3,1) {$(3,1)$};
    \end{axis}
\end{tikzpicture}
\end{minipage}

\begin{center}
The slope of the line is $-\frac{2}{3}$.    
\end{center}
\end{frame}


\begin{frame}
\textbf{\gls{distance vs. time graph}}: \glsdesc{distance vs. time graph}

\begin{center}
\begin{tikzpicture}
    \begin{axis}[height=6cm,width=7cm,
        axis lines=left,
        xlabel = {Time (hours)},
        ylabel = {Distance (miles)},
        ymin=0, ymax=80,
        xmin=0, xmax=5,
        ytick={0,20,...,80},
        xtick={0,1,...,5},
        minor y tick num=1,
        minor x tick num=1,
        grid=both,
        ]
        \addplot[very thick,domain=0:5] {15*x};
    \end{axis}
\end{tikzpicture}
\end{center}
\end{frame}

\begin{frame}
On a distance vs.\,time graph, the \textcolor{red}{\bfseries slope} of the line is the \textcolor{red}{\bfseries speed} of the object.

\begin{center}
\begin{tikzpicture}
    \begin{axis}[height=6cm,width=7cm,
        axis lines=left,
        xlabel = {Time (hours)},
        ylabel = {Distance (miles)},
        ymin=0, ymax=80,
        xmin=0, xmax=5,
        ytick={0,20,...,80},
        xtick={0,1,...,5},
        minor y tick num=1,
        minor x tick num=1,
        grid=both,
        ]
        \addplot[very thick,domain=0:5] {15*x};
    \end{axis}
\end{tikzpicture}
\end{center}
\end{frame}

\begin{frame}
\textit{Example.} Find the speed of the object below.

\begin{center}
\begin{tikzpicture}
    \begin{axis}[height=6cm,width=7cm,
        axis lines=left,
        xlabel = {Time (hours)},
        ylabel = {Distance (miles)},
        ymin=0, ymax=80,
        xmin=0, xmax=5,
        ytick={0,20,...,80},
        xtick={0,1,...,5},
        minor y tick num=1,
        minor x tick num=1,
        grid=both,
        clip=false
        ]
        \addplot[very thick,domain=0:5] {15*x};
        \addplot[red,only marks,mark=*] coordinates{(0,0)(4,60)};
        \node[red,shift={(0.5,20)}] at (0,0) {$(0,0)$};
        \node[red,shift={(-0.5,10)}] at (4,60) {$(4,60)$};
        \node[overlay,red,fill=black!8] at (4.5,25) {$\displaystyle m = \frac{60-0}{4-0} = 15$};
    \end{axis}
\end{tikzpicture}
\end{center}

The speed is 15 miles per hour.
\end{frame}

\section{Thu Aug 20 - area, speed vs time graph}

\begin{frame}
\textbf{\gls{area}}: \glsdesc{area}

\begin{center}
\begin{tikzpicture}
    \draw[thick,fill=black!10] (0,0) rectangle ++(3,2);
    \node[below] at (1.5,0) {$l$};
    \node[left] at (0,1) {$w$};
    \node at (1.3,-2) {$\displaystyle A_\text{\scriptsize rec} = lw$};
    \begin{scope}[shift={(5.2,0)}]
        \draw[thick,fill=black!10] (-1,0) -- (1,0) node[below,pos=0.5] {$b$} -- (0,2) -- cycle;
        \draw (0,0) -- (0,2) node[right,pos=0.4] {$h$};   
        \node at (0,-2) {$\displaystyle A_\text{\scriptsize tri} = \frac{1}{2}bh$};
    \end{scope}
    \begin{scope}[shift={(8.5,1)}]
        \draw[thick,fill=black!10] (0,0) circle (1);    
        \draw[thick] (0,0) -- (1,0) node[above,pos=0.5] {$r$};
        \node at (0,-3) {$\displaystyle A_\text{\scriptsize cir} = \pi r^2$};
    \end{scope}
\end{tikzpicture}
\end{center}
\end{frame}

% \begin{frame}
% The area of a rectangle is length times width.

% \begin{center}
% \begin{tikzpicture}[x=1.5cm,y=1.5cm]
%     \draw[thick,fill=black!10] (0,0) rectangle ++(3,2);
%     \node[below] at (1.5,0) {length $(l)$};
%     \node[left] at (0,1) {width $(w)$};
% \end{tikzpicture}

% \vspace{-1cm}

% {\large
% \begin{equation*}
%     \mathrm{area} = A = lw = \mathrm{length} \times \mathrm{width}
% \end{equation*}
% }
% \end{center}
% \end{frame}

\begin{frame}
The area of a rectangle is base times height.

\begin{center}
\begin{tikzpicture}[x=1.5cm,y=1.5cm]
    \draw[thick,fill=black!10] (0,0) rectangle ++(3,2);
    \node[below] at (1.5,0) {base $(b)$};
    \node[left] at (0,1) {height $(h)$};
\end{tikzpicture}

\vspace{-1cm}

{\large
\begin{equation*}
    \mathrm{area} = A = bh = \mathrm{base} \times \mathrm{height}
\end{equation*}
}
\end{center}
\end{frame}

\begin{frame}
\textit{Example.} Calculate the area of the rectangle below.

\begin{center}
\begin{tikzpicture}
    \draw[thick,fill=black!10] (0,0) rectangle ++(3,2);
    \node[below] at (1.5,0) {3};
    \node[left] at (0,1) {2};
\end{tikzpicture}
\end{center}

\begin{equation*}
    A = bh = (3)(2) = \boxed{6}
\end{equation*}
\end{frame}

\begin{frame}
\textbf{\gls{speed vs. time graph}}: \glsdesc{speed vs. time graph}

\begin{center}
\begin{tikzpicture}
    \begin{axis}[height=6cm,width=7cm,
        axis lines=left,
        xlabel = {Time (hours)},
        ylabel = {Speed (miles per hour)},
        ymin=0, ymax=20,
        xmin=0, xmax=5,
        ytick={0,5,...,20},
        xtick={0,1,...,5},
        minor y tick num=1,
        minor x tick num=1,
        grid=both,
        ]
        \addplot[very thick,domain=0:5] {15};
    \end{axis}
\end{tikzpicture}
\end{center}
\end{frame}

\begin{frame}
On a speed vs.\,time graph, the \alert<2>{area} under the line is the \alert<2>{distance} traveled by the object.

\begin{center}
\begin{tikzpicture}
    \begin{axis}[height=6cm,width=7.5cm,
        axis lines=left,
        xlabel = {Time (hours)},
        ylabel = {Speed (miles per hour)},
        ymin=0, ymax=20,
        xmin=0, xmax=5,
        ytick={0,5,...,20},
        xtick={0,1,...,5},
        minor y tick num=1,
        minor x tick num=1,
        grid=both,
        clip=false,
        ]
        \addplot[very thick,domain=0:5] {15};
        \begin{pgfonlayer}{background}
            \fill<2->[red!10] (0,0) rectangle (5,15);
        \end{pgfonlayer}
    \end{axis}
\end{tikzpicture}
\end{center}
\end{frame}

\begin{frame}
\textit{Example.} Find the total distance traveled in 4 hours.

\begin{center}
\begin{tikzpicture}
    \begin{axis}[height=6cm,width=7.5cm,
        axis lines=left,
        xlabel = {Time (hours)},
        ylabel = {Speed (miles per hour)},
        ymin=0, ymax=20,
        xmin=0, xmax=5,
        ytick={0,5,...,20},
        xtick={0,1,...,5},
        minor y tick num=1,
        minor x tick num=1,
        grid=both,
        clip=false,
        ]
        \addplot[very thick,domain=0:5] {15};
        \begin{pgfonlayer}{background}
            \fill[red!10] (0,0) rectangle (4,15);
        \end{pgfonlayer}
        \draw[very thick,red] (0,0) rectangle (4,15);
        \draw[|-|,red,thick] (0,16) -- (4,16) node[above,pos=0.5,fill=black!8] {4 hours};
        \draw[|-|,red,thick] (4.2,0) -- (4.2,15) node[right,pos=0.5,fill=black!8] {15 miles/hour};
        \node at (2,11) {$\displaystyle \mathrm{area} = \mathrm{base} \times \mathrm{height}$};
        \node at (2,6) {$\displaystyle \mathrm{area} = 4 \times 15 = 60$};
    \end{axis}
\end{tikzpicture}
\end{center}

The total distance traveled is 60 miles.
\end{frame}

\section{Fri Aug 21 - velocity and motion maps}

\begin{frame}{Warm-Up}
\begin{enumerate}
    \item What does position mean?
    \item The base SI unit more measuring position is\ldots
    \item Given an example of how to measure an object's position.
\end{enumerate}
\end{frame}

\againframe{5i5c0}

\begin{frame}{Warm-Up (cont.)}
    Displacement means\ldots

    \begin{enumerate}[A]
        \item \glsdesc{area}
        \item \glsdesc{distance}
        \item \alert<2>{\glsdesc{displacement}}
        \item \glsdesc{position}
    \end{enumerate}
\end{frame}



\begin{frame}
\textbf{instantaneous velocity}: \glsdesc{velocity1}
\vspace{1em}

velocity base SI units: meters per second (m/s)
\vspace{1em}

\textbf{\gls{velocity2}}: \glsdesc{velocity2}

\begin{center}
\begin{tikzpicture}[y=1cm,x=1cm]
    \foreach \x in {0,1,...,10} {
        \draw[ultra thin] (\x,-1mm) -- (\x,1mm) node[below=2mm] {$\x$};
    }
    \draw[->] (0,0) -- (10,0) node[below=6mm,pos=0.5] {Position (m)};
    \node[above] at (3,0) {\Strichmaxerl[2]};
    \draw[ultra thick,->] (2.6,0.5) -- (1.6,0.5) node[above,pos=0.5] {$v$};
\end{tikzpicture} 
\end{center}
\end{frame}

\begin{frame}
\textit{Example.} What is the direction of the object's velocity?

\begin{center}
\begin{tikzpicture}[y=1cm,x=1cm]
    \foreach \x in {0,1,...,10} {
        \draw[ultra thin] (\x,-1mm) -- (\x,1mm) node[below=2mm] {$\x$};
    }
    \draw[->] (0,0) -- (10,0) node[below=6mm,pos=0.5] {Position (m)};
    \node[above] at (3,0) {\Strichmaxerl[2]};
    \draw[ultra thick,->] (2.6,0.5) -- (1.6,0.5) node[above,pos=0.5] {$v$};
\end{tikzpicture} 
\end{center}

\pause 

{\color{red} \textit{Solution:} The velocity direction is negative, or leftward. Velocity is a vector, so it has magnitude and direction.}
\end{frame}

\begin{frame}
\textbf{\gls{average velocity}} ($\bar{v}$): \glsdesc{average velocity}
\vspace{-1em}

\begin{align*}
    \bar{v} = \frac{\Delta x}{\Delta t}
\end{align*}

\begin{center}
\def\arraystretch{1.5}
\begin{tabular}{|c|c||c|}
    \hline
    \thead{\bfseries Symbol} & \thead{\bfseries Meaning} & \thead{\bfseries Unit (symbol)} \\ \hline
    \uncover<2->{$\bar{v}$ & average velocity & {\small meters per second (m/s)}} \\ \hline
    \uncover<3->{$\Delta x$ & displacement & meters (m)} \\ \hline
    \uncover<4->{$\Delta t$ & total time & seconds (s)} \\ \hline
\end{tabular}
\end{center}
\end{frame}

\begin{frame}
\textit{Example.} The figure below tracks the one-dimensional motion of a particle. The total time for the trip is 3 seconds. What is the particle's average velocity?
\vspace{1em}

\begin{center}
\begin{tikzpicture}[y=3mm,x=8mm]
    \draw[black!10,xstep=1,ystep=2] (0,0) grid (12,4);
    \foreach \x in {0,1,...,12} {
        \draw[ultra thin] (\x,-1mm) -- (\x,1mm) node[below=3mm] {$\x$};
    }
    \draw[->] (0,0) -- (12,0) node[pos=0.5,below=2.5] {Position (m)};
    \draw[ultra thick,->,rounded corners=0.7mm] (11,1)
        -- (8,1) 
        -- (8,2) 
        -- (10,2) 
        -- (10,3) 
        -- (5,3);
    \node at (11.5,1) {$x_i$};
    \node at (4.5,3) {$x_f$};
\end{tikzpicture}
\end{center}

{\color{red} Step 1: Find displacement.\\ Step 2: Use average velocity formula ($\bar{v} = \Delta x/\Delta t$).}
\end{frame}

\begin{frame}
\vspace{-1em}
\begin{center}
\begin{tikzpicture}[y=3mm,x=8mm]
    \draw[black!10,xstep=1,ystep=2] (0,0) grid (12,4);
    \foreach \x in {0,1,...,12} {
        \draw[ultra thin] (\x,-1mm) -- (\x,1mm) node[below=3mm] {$\x$};
    }
    \draw[->] (0,0) -- (12,0) node[pos=0.5,below=2.5] {Position (m)};
    \draw[ultra thick,->,rounded corners=0.7mm] (11,1)
        -- (8,1) 
        -- (8,2) 
        -- (10,2) 
        -- (10,3) 
        -- (5,3);
    \node at (11.5,1) {$x_i$};
    \node at (4.5,3) {$x_f$};
\end{tikzpicture}
\end{center}

\begin{equation*}
    \Delta x = x_f - x_i = \SI{5}{m} - \SI{11}{m} = -\SI{6}{m}
\end{equation*}

\begin{equation*}
    \bar{v} = \frac{\Delta x}{\Delta t} = \frac{-\SI{6}{m}}{\SI{3}{s}} = \boxed{-\SI{2}{m/s}}
\end{equation*}

\end{frame}


\begin{frame}
\textbf{\gls{motion map}}: \uncover<2->{\glsdesc{motion map}}

\uncover<3->{
\begin{center}
\begin{tikzpicture}
    \draw[black!20,step=1] (0,0) grid (10,2);  
    \foreach \x in {2,3,...,7}{
        \fill ({2*(\x-2)},1) circle (0.9mm);
    }
\end{tikzpicture}
\end{center}
}
\end{frame}

\begin{frame}
\begin{center}
\begin{tikzpicture}[y=7.8mm,x=9.5mm]
    \draw[black!20,step=1] (0,0) grid (10,2);  
    \foreach \x in {2,3,...,7}{
        \fill<\x-> ({2*(\x-2)},1) circle (0.9mm);
    }
    \uncover<8->{
        \node[above] at (0,2) {$t = \SI{0}{s}$};
        \foreach \x in {1,2,...,5}{
            \node[above] at (2*\x,2) {\SI{\x}{s}};
        }
    }
    
    \draw[|-|] (0,-0.2) -- (1,-0.2) node[below,pos=0.5] {\SI{1}{m}};
    
    \uncover<9->{
        \draw[|-|] (0,3.2) -- ++(2,0) node[align=center,pos=0.5,above,font=\small] {1\textsuperscript{st} time\\[-1mm] interval};
    }
    \uncover<10->{
        \draw[-|] (2,3.2) -- ++(2,0) node[align=center,pos=0.5,above,font=\small] {2\textsuperscript{nd} time\\[-1mm] interval};
    }
    \uncover<11->{
        \draw[-|] (4,3.2) -- ++(2,0) node[align=center,pos=0.5,above,font=\small] {3\textsuperscript{rd} time\\[-1mm] interval};
    }
    \uncover<12->{
        \draw[-|] (6,3.2) -- ++(2,0) node[align=center,pos=0.5,above,font=\small] {4\textsuperscript{th} time\\[-1mm] interval};
    }
    \uncover<13->{
        \draw[-|] (8,3.2) -- ++(2,0) node[align=center,pos=0.5,above,font=\small] {5\textsuperscript{th} time\\[-1mm] interval};
    }
\end{tikzpicture}
\end{center}
\end{frame}

\section{Tue Aug 25 - position vs time, Lab: Dune Buggy}

\begin{frame}
\textbf{\gls{position vs. time graph}}: \glsdesc{position vs. time graph}

\begin{center}
\begin{tikzpicture}
    \begin{axis}[height=5cm,width=7cm,
        axis lines=left,
        ymin=0, ymax=6,
        xmin=0, xmax=10,
        ylabel = {Position (m)},
        xlabel = {Time (s)},
        grid = both,
        ytick={0,1,...,6},
        xtick={0,1,...,10},
    ]
    \addplot[black,very thick,domain=0:10] {0.5*x};
\end{axis}
\end{tikzpicture}
\end{center}
\end{frame}

\begin{frame}
\begin{center}
\begin{tikzpicture}
    \begin{axis}[height=5cm,width=7cm,
        axis lines=left,
        ymin=0, ymax=6,
        xmin=0, xmax=10,
        ylabel = {Position (m)},
        xlabel = {Time (s)},
        grid = both,
        ytick={0,1,...,6},
        xtick={0,1,...,10},
    ]
    \addplot[black,very thick,domain=0:10] {0.5*x};
    \addplot[black,very thick,densely dashed,domain=0:10] {5 - 0.4*x};
\end{axis}
\end{tikzpicture}
\end{center}

Lines on a \underline{position} vs time graph can have negative slopes. (Negative slopes NOT allowed for \underline{distance} vs time graphs.)
\end{frame}

\begin{frame}
\textit{Example.} What is the object's displacement between $t = \SI{2}{s}$ and $t = \SI{5}{s}$?

\begin{center}
\begin{tikzpicture}
    \begin{axis}[width=7cm,height=5cm,
        axis lines=left,
        ymin=0, ymax=10,
        xmin=0, xmax=10,
        ylabel = {Position (m)},
        xlabel = {Time (s)},
        grid = both,
        xtick={0,1,...,10},
        ytick={0,2,...,10},
        minor y tick num=1,
    ]
    \addplot[black,very thick]
        coordinates{(0,0)(4,8)(8,8)(10,4)};
\end{axis}
\end{tikzpicture}
\end{center}

{\color{red} \textit{Solution:}
\vspace{-1em}

\begin{equation*}
    \Delta x = x_f - x_i = \SI{8}{m} - \SI{4}{m} = \boxed{+\SI{4}{m}}
\end{equation*}
}
\end{frame}



\section{Thu Aug 27 - vectors and scalars}


\begin{frame}
\textbf{\gls{vector}}: \uncover<2->{\glsdesc{vector}}

\begin{center}
\begin{tikzpicture}
    \draw[black!12] (0,-1) grid (10,1);
    \draw<3->[line width=1mm,-{Triangle}] (2,0) -- (8,0);
    \node<4->[left=1mm,fill=black!10] at (2,0) {tail}; 
    \node<5->[right=1mm,fill=black!10] at (8,0) {head};
    \draw<6->[|-|,thick] (2,-1.2) -- (8,-1.2) node[below,pos=0.5] {magnitude};
\end{tikzpicture}
\end{center}
\end{frame}

\begin{frame}
\textit{Example.} If the length of each grid square is 1 unit, what is the magnitude of the vector?

\begin{center}
\begin{tikzpicture}
    \draw[black!12] (0,-1) grid (10,1);
    \draw[line width=1mm,-{Triangle}] (2,0) -- (8,0);
    \node[left=1mm,fill=black!10] at (2,0) {tail}; 
    \node[right=1mm,fill=black!10] at (8,0) {head};
    \draw<2->[|-|,thick] (2,-1.2) -- (8,-1.2) node[red,below,pos=0.5] {magnitude = 6 units} ;
\end{tikzpicture}
\end{center}

\uncover<2->{\color{red} \textit{Solution:} Counting squares, the magnitude is 6 units.}
\end{frame}



\begin{frame}
\textbf{\gls{head-to-tail method}}: \glsdesc{head-to-tail method}

\vspace{-5mm}

\begin{center}
\begin{tikzpicture}
    \draw[black!12] (0,0) grid (5,4);
    \uncover<2->{
        \draw[line width=1mm,-{Triangle}] (0,4) -- ++(2,0);
        \draw[line width=1mm,-{Triangle}] (5,3) -- ++(-3,0);
        \draw[line width=1mm,-{Triangle}] (0,2) -- ++(4,0);
        \draw[line width=1mm,-{Triangle}] (5,1) -- ++(-2,0);
        \draw[line width=1mm,-{Triangle}] (1,0) -- ++(3,0);
    }
\end{tikzpicture}%
\hspace{1mm}
\begin{tikzpicture}
    \draw[black!12] (0,0) grid (5,4);
    \node[overlay,below] at (0,0) {origin};
    \uncover<3->{
        \draw[line width=1mm,-{Triangle}] (0,0) -- ++(3,0);
    }
    \uncover<4->{
        \draw[line width=1mm,-{Triangle}] (3,1) -- ++(-2,0);
    }
    \uncover<5>{
        \draw[red,overlay] (3,0.5) circle (0.7);
        \node[red] at (3,1.5) {head-to-tail};
    }
    \uncover<6->{
        \draw[line width=1mm,-{Triangle}] (1,2) -- ++(4,0);
    }
    \uncover<7>{
        \draw[red,overlay] (1,1.5) circle (0.7);
        \node[red] at (1,2.6) {head-to-tail};
    }
    \uncover<8->{
        \draw[line width=1mm,-{Triangle}] (5,3) -- ++(-3,0);
    }
    \uncover<9>{
        \draw[red,overlay] (4.7,2.5) circle (0.7) node[above=0.7] {h-2-t};
    }
    \uncover<10->{
        \draw[line width=1mm,-{Triangle}] (2,4) -- ++(2,0);
    }
    \uncover<11>{
        \draw[red,overlay] (2,3.5) circle (0.7) node[above=0.7] {h-2-t};
    }
    \uncover<12>{
        \node[overlay,red] at (4.2,4.4) {end};
    }
    \uncover<13>{
    \draw[line width=1mm,-{Triangle},red] (0,4.6) -- ++(4,0) node[pos=0.5,above=0.5mm] {vector sum};
    }
\end{tikzpicture}
\end{center}
\end{frame}

\begin{frame}
\textbf{\gls{vector sum}}: \glsdesc{vector sum}
\vspace{1em}

\begin{center}
    vector sum = 1 vector that combines all vectors
\end{center}
\end{frame}




\end{document}

